\section{Exact Closed-Horndeski Kernel for the First Physical Scalar Mode}

The previous section uses the Bellini--Sawicki functions to organize the cubic-Galileon background. For the curvature-scale mode, however, the flat-FLRW perturbation formulas are not sufficient. We therefore specialize the non-flat Horndeski quadratic action derived by Akama and Kobayashi \citep{AkamaKobayashi2019}.

For
\[
G_2=X-V(T),\qquad
G_3=\frac{2X}{\Lambda^3},\qquad
G_4=\frac{M_{\rm Pl}^2}{2},\qquad
G_5=0,
\]
their tensor coefficients reduce exactly to
\[
\mathcal F_T=\mathcal G_T=M_{\rm Pl}^2.
\]
Because \(G_4\) is constant and \(G_5=0\), the explicit curvature contributions to their \(\Theta\) and \(\Sigma\) vanish. The remaining coefficients are
\[
\boxed{
\Theta
=
HM_{\rm Pl}^2-\frac{\dot T^3}{\Lambda^3}
}
\]
and
\[
\boxed{
\Sigma
=
\frac12\dot T^2
+\frac{12H\dot T^3}{\Lambda^3}
-3M_{\rm Pl}^2H^2.
}
\]
Equivalently,
\[
\Theta
=
HM_{\rm Pl}^2
\left(1-\frac{\alpha_B}{2}\right).
\]

Define
\[
r(t)
\equiv
\frac{M_{\rm Pl}^2\Sigma}{\Theta^2},
\qquad
\kappa(t)\equiv\frac{K}{a^2}.
\]
The first physical scalar harmonic has alternative-convention index \(N=3\), so
\[
k^2=K(N^2-1)=8K.
\]
The two Akama--Kobayashi scalar constraints then reduce to
\[
\Sigma\,\delta n
+8\Theta\kappa\,\chi
+3\Theta\dot\zeta
+5M_{\rm Pl}^2\kappa\,\zeta=0,
\]
\[
\Theta\,\delta n
-M_{\rm Pl}^2\dot\zeta
-M_{\rm Pl}^2\kappa\,\chi=0.
\]
They can be solved algebraically:
\[
\boxed{
\delta n
=
\frac{5M_{\rm Pl}^2}{\Theta(r+8)}
\left[
\dot\zeta
-\frac{M_{\rm Pl}^2\kappa}{\Theta}\zeta
\right]
}
\]
and
\[
\boxed{
\chi
=
-\frac{r+3}{\kappa(r+8)}\dot\zeta
-\frac{5M_{\rm Pl}^2}{\Theta(r+8)}\zeta.
}
\]

The reduced kinetic coefficient of this \(n=2\) mode is therefore
\[
\boxed{
\mathcal G_{S,2}
=
M_{\rm Pl}^2
\frac{5(r+3)}{r+8}.
}
\]
Using the cubic-Galileon \(\alpha\)-functions gives the exact identity
\[
r+3
=
\frac{D_{\rm kin}}
{2(1-\alpha_B/2)^2},
\]
so that
\[
\boxed{
\mathcal G_{S,2}
=
M_{\rm Pl}^2
\frac{5D_{\rm kin}}
{D_{\rm kin}
+10(1-\alpha_B/2)^2}.
}
\]
Thus, provided \(\Theta\neq0\), positivity of the previously computed \(D_{\rm kin}\) is sufficient for positivity of the pure gravity--scalar kinetic coefficient of the first physical closed-universe scalar harmonic.

For completeness, the corresponding non-flat Horndeski gradient coefficient is
\[
\mathcal F_{S,2}
=
\frac1a\frac{d}{dt}
\left[
\frac{5}{r+8}
\left(
\frac{aM_{\rm Pl}^4}{\Theta}
\right)
\right]
-M_{\rm Pl}^2
+
\frac{5}{r+8}
\frac{M_{\rm Pl}^6}{\Theta^2}
\frac{K}{a^2}.
\]
The repository evaluates this expression as a gravity--scalar diagnostic only. The late-time R1 background contains matter and radiation, whose perturbations introduce additional dynamical variables and source terms. Consequently, \(\mathcal F_{S,2}\) is not promoted here to a complete matter-era stability or light-cone transfer result.

This section closes an important part of the bridge: the first physical harmonic has an explicit constrained scalar--metric kernel rather than only a geometric identification. The remaining observable calculation is the coupled matter/radiation plus topographic-response system and its gauge-invariant transformation to the metric and velocity variables entering distance, growth, and lensing observables.
